\Introduction

\paragraphe{}{
Les documents occupent une place essentielle en informatique appliquée. Pourtant, ils continuent d'être négligés et délaissés. Pour plusieurs, cela peut paraître anodin, mais les répercussions que peut engendrer une absence ou une mauvaise documentation sont réelles. Ainsi, la mauvaise documentation a été pointée comme l'un des facteurs déterminants dans les tragédies causées en 2018 et 2019 par le logiciel \emphase{Maneuvering Characteristics Augmentation System} de Boeing \cite{johnston_boeing_2019}. Plus de trente ans auparavant, la mauvaise documentation avait été aussi désignée comme un facteur important dans la tragédie du Therac-25, qui deviendra par la suite un cas d'école \cite{leveson_therac_2017, leveson_investigation_1993}. Malgré ces exemples, et de très nombreux autres, les pratiques n’accordent généralement pas l’importance et les moyens qu’elles devraient à la documentation, et ce, même dans les secteurs critiques\footnote{Le terme \emphase{secteur critique} signifie un secteur d'infrastructure critique ou essentielle à la société (énergie, transport, santé, défense, télécommunication, etc.) [Gouvernement du Canada, site web, «~Infrastructures essentielles (IE) du Canada~», \url{https://www.securitepublique.gc.ca/cnt/ntnl-scrt/crtcl-nfrstrctr/cci-iec-fr.aspx} (consulté le 2025-01-24)].}.
}{}{}

\paragraphe{}{
Exposer les mécanismes existants entre les documents et les logiciels ainsi que leurs rôles devrait suffire à convaincre de l'importance de la documentation ceux qui pratiquent l'informatique appliquée, puisqu'il s'agit de l'application d'une science. Pourtant, les années s'écoulent et la vaste majorité des méthodes relatives aux documents en informatique appliquée continuent à être négligées. Par conséquent, la qualité des artefacts informatiques en pâtit et les incidents, bien qu'évitables, se poursuivent. Certaines approches ont été proposées afin d'éviter de reproduire les mêmes erreurs, mais, globalement, le cycle ne se brise pas. Ce mémoire et la preuve de concept associée découlent de la volonté de contribuer à remédier à cette situation. Les chapitres compris dans ce mémoire sont~: la revue de littérature, une définition du problème, la réalisation d'une solution et l'évaluation de la solution. 
}{}{}

\paragraphe{}{
En premier lieu, le chapitre 1 comprend une revue de littérature qui s’efforce de cerner la nature et l’importance des relations entre les documents, les modèles et les connaissances. Le chapitre est divisé en trois sections et chacune de ses sections en trois sous-sections : les objectifs, les écueils à ces objectifs et les approches pour surmonter ces écueils. Ce chapitre sert de fondation aux autres chapitres. Cette revue de littérature montre, entre autres, que l'atteinte d'un niveau de qualité pour des logiciels est indissociable de la qualité des modèles utilisés. Les documents en sont le véhicule, ils permettent de faire la validation, la réalisation et l'établissement du logiciel. Ainsi, des documents erronés conduisent à des modèles erronés et des modèles erronés conduisent à des logiciels erronés. 
}{}{}

\paragraphe{}{
Ensuite, le chapitre 2 présente une définition du problème portant sur la formalisation et la compréhension des documents servant aux modèles. Cette définition intègre plusieurs approches qui se sont avérées efficaces sans pour autant être largement utilisées. Après avoir présenté le problème, une description d'un modèle du problème est présentée ainsi que les attentes et les besoins.
}{}{}

\paragraphe{}{
Puis, le chapitre 3 présente la réalisation d'une solution à ce problème. Cela débute avec des exigences répondant au problème. Ensuite, une conception est décrite et de nouvelles exigences apparaissent. S'ensuivent des explications sur l'implémentation et sur les instruments utilisés. Le chapitre se termine par un tableau identifiant les besoins comblés par les exigences ainsi qu'un prolongement à la solution développée.}{}{}

\paragraphe{}{Par la suite, le chapitre~4 présente l'utilisation de la solution. Cela commence en exposant des situations où la solution pourrait être efficace. Ensuite, des explications sont données sur l'utilisation de la solution à travers un exemple.}{}{}

\paragraphe{}{
Après, le chapitre~5 présente l'évaluation de la solution découpée en deux étapes~: la première consiste en la vérification des exigences, à l'aide de cas de tests, et la deuxième, en une proposition de validation de certaines attentes. Cette proposition de validation comporte la description partielle d'un protocole expérimental. Ce découpage est aussi celui des caractéristiques de qualité. Les caractéristiques de qualité absolues se vérifiant par des cas de tests, tandis que celles relatives se validant par la proposition de validation.
}{}{}


\paragraphe{}{
Finalement, la conclusion présente les différentes contributions de ce mémoire ainsi que trois prolongements. Les contributions comprennent une synthèse des thèmes de la documentation, de la modélisation et des connaissances dans le domaine de l'informatique appliquée, ainsi que la démarche ayant mené à une solution prenant la forme d'une preuve de concept. Les prolongements s'inscrivent dans la continuité de cette preuve de concept~: les deux premiers portent sur l'intégration et la comparaison des méthodes formelles et de l'intelligence artificielle, et le troisième vise à rendre cette preuve de concept plus générique encore.
}{}{}